\documentclass[10pt,a4paper]{amsart}
\usepackage[margin=19mm]{geometry}
\usepackage{amsmath,amssymb,mathtools}
\usepackage[scr=rsfs,cal=euler]{mathalfa}
\usepackage{xfrac}
\usepackage[unicode,hidelinks,bookmarks=false]{hyperref}
\newcommand{\ie}{i.e.\ }
\newcommand{\cf}{cf.\ }
\newcommand{\resp}{respectively\ }
\newcommand{\Subsection}{\subsection}
\providecommand{\nobreakdash}{\nobreak\hbox{-}}
\setlength{\emergencystretch}{2em}
\multlinegap0pt
\usepackage{bm}
\usepackage{amssymb}
\newtheorem{lemma}{Lemma}
\renewcommand{\P}{\mathbb{P}}
\newcommand{\T}{\mathbb{T}}
\newcommand{\X}{\bm{X}}
\newcommand{\Y}{\bm{Y}}
\newcommand{\abs}[1]{\left\lvert #1 \right\rvert}
\newcommand{\mc}[1]{\mathcal{#1}}
\newcommand{\norm}[1]{\left\lVert #1 \right\rVert}
\newcommand{\p}{\partial}
\title{A nonlocal curve evolution for an immersed elastic filament: global existence and convergence to resistive force theory}
\author{Laurel Ohm}
\date{Source: arXiv:2604.09496v1 (CC BY 4.0)}
\begin{document}
\maketitle

\noindent\emph{Source and licence.} A nonlocal curve evolution for an immersed elastic filament: global existence and convergence to resistive force theory. Version arXiv:2604.09496v1 (CC BY 4.0), distributed by arXiv under the Creative Commons Attribution 4.0 International licence, http://creativecommons.org/licenses/by/4.0/, as recorded in the arXiv licence metadata for this exact version and shown on the abstract page https://arxiv.org/abs/2604.09496v1. Full licence notice: \texttt{licenses/arxiv-2604.09496-CC-BY-4.0.txt}.\par\smallskip

\noindent\textit{Editorial scope.} Counted unit: the article's lemma lem:tensionXY\_diff (source0.tex lines 1797--1808). The article's complete original proof of it is delivered with it (source0.tex lines 1811--2009, one \texttt{proof} environment of 199 lines). No mathematics is written, repaired or re-derived: the statement and the proof are the article's own lines, in the article's own notation, and no retained line is substituted or re-flowed.  Retained uncounted as the source-local context the proof needs: lines 518--522; lines 525--533; lines 1633--1643; lines 1653--1664; lines 1775--1780.  Nothing was omitted from any retained span, so the package carries no omission record.  No retained line contains a citation, so the wrapper carries no bibliography and no bibliography entry.  No retained line contains a figure environment, an image-inclusion command or drawing code, so no image is delivered and the package declares no asset dependency.  Editorial formatting only, in addition: an AMS \texttt{amsart} preamble that restates the article's own declarations and macros used by the retained lines, namely the environments \texttt{lemma} on separate counters, together with the standard packages those lines need.  The retained environments print in this document as Lemma 1 in order of appearance, whereas the article prints the same items with the numbering of its own full document.  The complete native source of the article as distributed in the arXiv e-print for this exact version is bundled unchanged as \texttt{sources/198212/source0.tex}.
\begin{equation}\label{eq:inex_ids}
  \X_s\cdot\X_{ss}=0\,, \quad 
  \abs{\X_{ss}}^2 =-\X_s\cdot\X_{sss}\,, \quad
  3\X_{sss}\cdot\X_{ss} = -\X_s\cdot\X_{ssss}\,.
\end{equation}

\begin{equation}\label{eq:interps}
\begin{aligned}
  \norm{h}_{\dot H^{1/2}}&\le c\norm{\p_sh}_{L^2}^{1/2}\norm{h}_{L^2}^{1/2}\,,\\
%
  \norm{h_s}_{L^2}&\le c\norm{h_s}_{\dot H^{1/2}}^{2/3}\norm{h}_{L^2}^{1/3}\,,\\
  %
  \norm{h}_{L^\infty} &\le c\norm{h_s}_{\dot H^{1/2}}^{1/3}\norm{h}_{L^2}^{2/3}\,.
\end{aligned}
\end{equation}

\begin{lemma}[Inverse multiplier operators]\label{lem:Tinverse}
  For $\theta\ge 0$, the inverse operators 
  \begin{equation}
    T_{m_\epsilon^{\rm t}}^{-\theta}(\cdot) = \mc{F}^{-1}\big[m_\epsilon^{\rm t}(k)^{-\theta}\mc{F}[\cdot]\big]\,, \quad
    T_{m_\epsilon^{\rm n}}^{-\theta}(\cdot) = \mc{F}^{-1}\big[m_\epsilon^{\rm n}(k)^{-\theta}\mc{F}[\cdot]\big]
  \end{equation}
satisfy, for $j={\rm t},{\rm n}$,
  \begin{equation}
    \norm{T_{m_\epsilon^j}^{-\theta}\bm{f}}_{H^\ell} \le c\abs{\log\epsilon}^{-\theta}\norm{\bm{f}}_{H^{\ell+\theta}}\,.
  \end{equation}
\end{lemma}

\begin{lemma}[Multiplier operator and RFT difference]\label{lem:mult_RFT_diff}
Given $\bm{\varphi}\in H^{-1/2}(\T)$ and $\bm{f}\in L^2(\T)$, we may estimate 
 \begin{equation}
  \begin{aligned}
  \abs{\int_\T\bm{\varphi}\cdot\bigg(T_{m_\epsilon^{\rm t}}-\frac{\abs{\log\epsilon}}{2\pi}\bigg)[\bm{f}]\,ds} 
  &\le c\norm{T_{m_\epsilon^{\rm t}}^{1/2}\bm{\varphi}}_{L^2}\norm{\bm{f}}_{L^2}\\
  %
   \abs{\int_\T\bm{\varphi}\cdot\bigg(T_{m_\epsilon^{\rm n}}-\frac{\abs{\log\epsilon}}{4\pi}\bigg)[\bm{f}]\,ds} 
   &\le c\norm{T_{m_\epsilon^{\rm n}}^{1/2}\bm{\varphi}}_{L^2}\norm{\bm{f}}_{L^2}\,.
\end{aligned}
\end{equation}
\end{lemma}

\begin{equation}\label{eq:tauYbd}
\begin{aligned}
  \norm{\tau_Y}_{H^1} &\le c\norm{\Y_{ss}}_{H^1}^2 \,, \\
  \norm{\tau_Y}_{H^2}&\le c(\norm{\Y_s}_{H^2}^4 + \norm{\Y_s}_{H^2}\norm{\Y_s}_{H^3})\,.
\end{aligned}
\end{equation}

\begin{lemma}[Tension difference]\label{lem:tensionXY_diff}
  The difference $\tau_W=\tau_X-\tau_Y$ between the tensions corresponding to curves $\X$ and $\Y$, respectively, satisfies the bound 
\begin{equation}\label{eq:tauWbd}
\begin{aligned}
  \norm{T_{m_\epsilon^{\rm t}}^{1/2}((\tau_W)_s\X_s)}_{L^2} &+ 
  \norm{T_{m_\epsilon^{\rm n}}^{1/2}(\tau_W\X_{ss})}_{L^2}
  \le c\bigg(\abs{\log\epsilon}^{1/2}\norm{\bm{W}_s}_{H^1}(1+\norm{\Y_s}_{H^3}^2)\\
  & + \bigg(\norm{T_{m_\epsilon^{\rm t}}^{1/2}\bm{W}_{ssss}}_{L^2}+ 1+\norm{\Y_s}_{H^3}^2\bigg)\big(1+\norm{\X_s}_{H^1}+\norm{\Y_s}_{H^3}^2\big)\bigg)\,,
\end{aligned}
\end{equation}
where $\bm{W}=\X-\Y$.
\end{lemma}

\begin{proof}
For any $\varphi\in H^{1/2}(\T)$, defining the bilinear forms
\begin{equation}
\begin{aligned}
  \mc{B}_\epsilon^X(\tau_X,\varphi) &= \int_\T \overline{\mc{L}}_\epsilon(\X)[(\tau_X\X_s)_s] \cdot(\varphi\X_s)_s\,ds \,,\\
  \mc{B}_\epsilon^Y(\tau_Y,\varphi) &= \frac{\abs{\log\epsilon}}{4\pi}\int_\T \big(({\bf I}+\Y_s\otimes\Y_s)(\tau_Y\Y_s)_s\big) \cdot(\varphi\Y_s)_s\,ds\,,
\end{aligned}
\end{equation}
we may write the difference between the weak form of the equations satisfied by $\tau_X$ and $\tau_Y$ as 
\begin{equation}\label{eq:diffXY_weak}
\begin{aligned}
  &\mc{B}_\epsilon^X(\tau_X,\varphi) - \mc{B}_\epsilon^Y(\tau_Y,\varphi) \\
  &\qquad = \underbrace{\int_\T \bigg(\overline{\mc{L}}_\epsilon(\X)[\X_{ssss}]\cdot(\varphi\X_s)_s - \frac{\abs{\log\epsilon}}{4\pi}\big(({\bf I}+\Y_s\otimes\Y_s)\Y_{ssss}\big)\cdot(\varphi\Y_s)_s\bigg)\,ds }_{{\rm RHS}}\,.
\end{aligned}
\end{equation}


The left hand side of \eqref{eq:diffXY_weak} may be written as 
\begin{equation}
\begin{aligned}
  \mc{B}_\epsilon^X(\tau_X,\varphi) &- \mc{B}_\epsilon^Y(\tau_Y,\varphi)= B^\parallel +B^\perp \,,\\
  B^\parallel &= \int_\T \bigg(\big(\P_{\X_s}T_{m_\epsilon^{\rm t}}((\tau_X)_s\X_s)\big)\cdot(\varphi\X_s)_s -\frac{\abs{\log\epsilon}}{2\pi}(\tau_Y)_s\Y_s\cdot(\varphi\Y_s)_s \bigg)\,ds \\
  B^\perp &= \int_\T \bigg( \big(\P_{\X_s}^\perp T_{m_\epsilon^{\rm n}}(\tau_X\X_{ss})\big)\cdot(\varphi\X_s)_s  -\frac{\abs{\log\epsilon}}{4\pi}\tau_Y\Y_{ss} \cdot(\varphi\Y_s)_s \bigg)\,ds\,.
\end{aligned}
\end{equation}
%
We may further write 
\begin{equation}
\begin{aligned}
  B^\parallel &= \int_\T \bigg(\big(\X_s\cdot T_{m_\epsilon^{\rm t}}((\tau_X)_s\X_s)\big) -\frac{\abs{\log\epsilon}}{2\pi}(\tau_Y)_s \bigg)\varphi_s\,ds \\
  &= \int_\T \bigg(\big(\X_s\cdot T_{m_\epsilon^{\rm t}}((\tau_W)_s\X_s)\big) + \X_s\cdot\bigg(T_{m_\epsilon^{\rm t}}-\frac{\abs{\log\epsilon}}{2\pi}\bigg)\big[(\tau_Y)_s\X_s\big] \bigg)\varphi_s\,ds
\end{aligned}
\end{equation}
as well as 
\begin{equation}
\begin{aligned}
 B^\perp &= \int_\T \bigg( \X_{ss}\cdot T_{m_\epsilon^{\rm n}}(\tau_X\X_{ss})  -\frac{\abs{\log\epsilon}}{4\pi}\tau_Y\abs{\Y_{ss}}^2 \bigg)\varphi\,ds \\
  &= \int_\T \bigg( \X_{ss}\cdot T_{m_\epsilon^{\rm n}}(\tau_W\X_{ss}) + \X_{ss}\cdot\bigg(T_{m_\epsilon^{\rm n}} -\frac{\abs{\log\epsilon}}{4\pi}\bigg)[\tau_Y\X_{ss}] \\
  &\qquad\quad + \frac{\abs{\log\epsilon}}{4\pi}\tau_Y(\abs{\X_{ss}}^2-\abs{\Y_{ss}}^2) \bigg)\varphi\,ds \,.
\end{aligned}
\end{equation}
%
We may thus rewrite the left hand side of \eqref{eq:diffXY_weak} as 
\begin{equation}
\begin{aligned}
  \mc{B}_\epsilon^X(\tau_X,\varphi) &- \mc{B}_\epsilon^Y(\tau_Y,\varphi)= \mc{B}_\epsilon^X(\tau_W,\varphi) + {\rm LHS}\,,\\
  {\rm LHS}&= \int_\T \bigg( \varphi_s\X_s\cdot\bigg(T_{m_\epsilon^{\rm t}}-\frac{\abs{\log\epsilon}}{2\pi}\bigg)\big[(\tau_Y)_s\X_s\big] \\
  &\quad+ \varphi\X_{ss}\cdot\bigg(T_{m_\epsilon^{\rm n}} -\frac{\abs{\log\epsilon}}{4\pi}\bigg)[\tau_Y\X_{ss}]  + \frac{\abs{\log\epsilon}}{4\pi}\tau_Y(\abs{\X_{ss}}^2-\abs{\Y_{ss}}^2)\varphi \bigg)\,ds \,.
\end{aligned}
\end{equation}
%
%
Using Lemma \ref{lem:mult_RFT_diff} and the $\tau_Y$ bound \eqref{eq:tauYbd}, we may bound the terms appearing in LHS as 
\begin{equation}\label{eq:LHS_diffbd}
\begin{aligned}
  \abs{{\rm LHS}} &\le c\bigg( \norm{T_{m_\epsilon^{\rm t}}^{1/2}(\varphi_s\X_s)}_{L^2}\norm{(\tau_Y)_s\X_s}_{L^2} + 
  \norm{T_{m_\epsilon^{\rm n}}^{1/2}(\varphi\X_{ss})}_{L^2}\norm{\tau_Y\X_{ss}}_{L^2}\\
  &\qquad + \abs{\log\epsilon}\norm{\bm{W}_s}_{H^1}\norm{\tau_Y}_{L^\infty}(\norm{\bm{W}_{ss}}_{L^\infty}+\norm{\Y_{ss}}_{L^\infty})\norm{\varphi}_{L^2} \bigg) \\
  %
  %
  &\le c\bigg( \norm{T_{m_\epsilon^{\rm t}}^{1/2}(\varphi_s\X_s)}_{L^2}\norm{\Y_{ss}}_{H^1}^2 + 
  \norm{T_{m_\epsilon^{\rm n}}^{1/2}(\varphi\X_{ss})}_{L^2}\norm{\Y_{ss}}_{H^1}^2\norm{\X_s}_{H^1}\\
  &\qquad + \abs{\log\epsilon}\norm{\bm{W}_s}_{H^1}\norm{\Y_{ss}}_{H^1}^2(\norm{\bm{W}_{sss}}_{\dot H^{1/2}}^{2/3}\norm{\bm{W}_s}_{H^1}^{1/3}+\norm{\Y_{ss}}_{H^1})\times\\
  &\qquad\qquad\times\norm{T_{m_\epsilon^{\rm n}}^{-1}T_{m_\epsilon^{\rm n}}(\varphi \X_s)\cdot\X_s}_{L^2} \bigg)  \\
  %
  %
  &\le c\norm{\Y_{ss}}_{H^1}^2\bigg( \norm{T_{m_\epsilon^{\rm t}}^{1/2}(\varphi_s\X_s)}_{L^2} + 
  \norm{T_{m_\epsilon^{\rm n}}^{1/2}(\varphi\X_{ss})}_{L^2}\norm{\X_s}_{H^1}\\
  &\qquad + \norm{\bm{W}_s}_{H^1}(\norm{\bm{W}_{sss}}_{\dot H^{1/2}}^{2/3}\norm{\bm{W}_s}_{H^1}^{1/3}+\norm{\Y_{ss}}_{H^1})\norm{T_{m_\epsilon^{\rm n}}(\varphi \X_s)_s}_{L^2} \bigg) \\
  %
  %
  &\le c\bigg( \norm{T_{m_\epsilon^{\rm t}}^{1/2}(\varphi_s\X_s)}_{L^2} + 
  \norm{T_{m_\epsilon^{\rm n}}^{1/2}(\varphi\X_{ss})}_{L^2}\bigg)\norm{\Y_{ss}}_{H^1}^2\big(1 + \norm{\X_s}_{H^1}^2 + \norm{\Y_{ss}}_{H^1}^2 \\
  &\qquad\qquad + \abs{\log\epsilon}^{1/2}\norm{\bm{W}_s}_{H^1}^{4/3}\norm{\bm{W}_{sss}}_{\dot H^{1/2}}^{2/3}\big)\,.
\end{aligned}
\end{equation}
Here we have also exploited Lemma \ref{lem:Tinverse} to bound $T_{m_\epsilon^{\rm n}}^{-1}$, and have used that we may opt not to use the regularity gain from $T_{m_\epsilon^{\rm n}}^{1/2}$ to avoid an additional factor of $\epsilon^{-1/2}$ from the high wavenumber scaling of $m_\epsilon^{\rm n}(k)$. We have further used that $m_\epsilon^{\rm n}(k)\lesssim m_\epsilon^{\rm t}(k)\lesssim m_\epsilon^{\rm n}(k)$, so we may interchange $T_{m_\epsilon^{\rm t}}$ and $T_{m_\epsilon^{\rm n}}$ in the above bounds. 



We next turn to the right hand side terms RHS of \eqref{eq:diffXY_weak}, which we may write as 
\begin{equation}
\begin{aligned}
  {\rm RHS} &= R^\parallel + R^\perp \,,\\
  R^\parallel &= \int_\T \bigg(\X_s\cdot T_{m_\epsilon^{\rm t}}(\X_s\cdot\X_{ssss}\X_s) - \frac{\abs{\log\epsilon}}{2\pi}\Y_s\cdot\Y_{ssss}\bigg)\varphi_s\,ds \\
  %
  R^\perp &= \int_\T \bigg(\X_{ss}\cdot T_{m_\epsilon^{\rm n}}(\P_{\X_s}^\perp\X_{ssss}) - \frac{\abs{\log\epsilon}}{4\pi}\Y_{ssss}\cdot\Y_{ss} \bigg)\varphi\,ds \,.
\end{aligned}
\end{equation}
%
%
Using the inextensibility identities \eqref{eq:inex_ids}, we may write
\begin{equation}
\begin{aligned}
  R^\parallel &= -3\int_\T \varphi_s\X_s\cdot T_{m_\epsilon^{\rm t}}\big( (\X_{ss}\cdot\bm{W}_{sss}+\bm{W}_{ss}\cdot\Y_{sss})\X_s\big) \,ds\\
  &\qquad -3\int_\T \varphi_s\X_s\cdot \bigg(T_{m_\epsilon^{\rm t}}- \frac{\abs{\log\epsilon}}{2\pi}\bigg)\big[\X_s \Y_{ss}\cdot\Y_{sss}\big]\,ds\\
  %
  &= -3\int_\T \varphi_s\X_s\cdot T_{m_\epsilon^{\rm t}}\bigg( \frac{1}{2}\p_s\abs{\bm{W}_{ss}}^2\X_s+\p_s(\Y_{ss}\cdot\bm{W}_{ss})\X_s\bigg) \,ds\\
  &\qquad -\frac{3}{2}\int_\T \varphi_s\X_s\cdot \bigg(T_{m_\epsilon^{\rm t}}- \frac{\abs{\log\epsilon}}{2\pi}\bigg)\big[\p_s\abs{\Y_{ss}}^2\X_s\big]\,ds\,.
\end{aligned}
\end{equation}
Using Lemma \ref{lem:mult_RFT_diff} and the interpolation inequalities \eqref{eq:interps}, we may estimate $R^\parallel$ as 
\begin{equation}\label{eq:Rpar_bd}
\begin{aligned}
  \abs{R^\parallel} &\le c \norm{T_{m_\epsilon^{\rm t}}^{1/2}(\varphi_s\X_s)}_{L^2}
  \bigg(\norm{T_{m_\epsilon^{\rm t}}^{1/2}(\p_s\abs{\bm{W}_{ss}}^2\X_s)}_{L^2}+\norm{T_{m_\epsilon^{\rm t}}^{1/2}(\p_s(\bm{W}_{ss}\cdot\Y_{ss})\X_s)}_{L^2}\\
  &\hspace{5cm} + \norm{\X_s \Y_{ss}\cdot\Y_{sss}}_{L^2}\bigg)\\
  %
  &\le c\norm{T_{m_\epsilon^{\rm t}}^{1/2}(\varphi_s\X_s)}_{L^2}
  \bigg(\abs{\log\epsilon}^{1/2}\norm{\abs{\bm{W}_{ss}}^2}_{\dot H^1}
  +\abs{\log\epsilon}^{1/2}\norm{\bm{W}_{ss}\cdot\Y_{ss}}_{\dot H^1}+ \norm{\abs{\Y_{ss}}^2}_{\dot H^1}\bigg)\\
    %
  &\le c\norm{T_{m_\epsilon^{\rm t}}^{1/2}(\varphi_s\X_s)}_{L^2}
  \bigg(\abs{\log\epsilon}^{1/2}\norm{\bm{W}_{sss}}_{\dot H^{1/2}}(\norm{\bm{W}_s}_{H^1}+\norm{\Y_s}_{H^3})\\
  &\hspace{4cm}
  +\abs{\log\epsilon}^{1/2}\norm{\bm{W}_s}_{H^1}\norm{\Y_s}_{H^3}
  + \norm{\Y_{ss}}_{H^1}^2\bigg)\\
  %
  &\le c\norm{T_{m_\epsilon^{\rm t}}^{1/2}(\varphi_s\X_s)}_{L^2}
  \bigg(\norm{T_{m_\epsilon^{\rm t}}^{1/2}\bm{W}_{ssss}}_{L^2}(\norm{\bm{W}_s}_{H^1}+\norm{\Y_s}_{H^3}) + \norm{\Y_{ss}}_{H^1}^2\\
  &\hspace{4cm}
  +\abs{\log\epsilon}^{1/2}\norm{\bm{W}_s}_{H^1}\norm{\Y_s}_{H^3}
  \bigg)\,.
\end{aligned}
\end{equation}
Here, as in the estimate \eqref{eq:LHS_diffbd} for LHS, we have opted to forego the regularity gain from $T_{m_\epsilon^{\rm t}}^{1/2}$ in order to avoid an additional factor of $\epsilon^{-1/2}$ from the high wavenumber behavior of $m_\epsilon^{\rm t}(k)$. Furthermore, in the last line, we have used that 
\begin{equation}
\norm{\bm{W}_{sss}}_{\dot H^{1/2}}=\|T_{m_\epsilon^{\rm t}}^{-1/2}T_{m_\epsilon^{\rm t}}^{1/2}\bm{W}_{sss}\|_{\dot H^{1/2}}\le c\abs{\log\epsilon}^{-1/2}\|T_{m_\epsilon^{\rm t}}^{1/2}\bm{W}_{sss}\|_{\dot H^1}\,,
\end{equation}
by Lemma \ref{lem:Tinverse}. 

We may similarly write $R^\perp$ as 
\begin{equation}
\begin{aligned}
  R^\perp &= \int_\T \bigg(\X_{ss}\cdot T_{m_\epsilon^{\rm n}}(\X_{ssss}+3\X_{ss}\cdot\X_{sss}\X_s) - \frac{\abs{\log\epsilon}}{4\pi}(\Y_{ssss}+3\Y_{ss}\cdot\Y_{sss}\Y_s)\cdot\Y_{ss} \bigg)\varphi\,ds \\
  %
  &= \int_\T \varphi\X_{ss}\cdot T_{m_\epsilon^{\rm n}}\big(\bm{W}_{ssss}+3\X_{ss}\cdot\bm{W}_{sss}\X_s +3\bm{W}_{ss}\cdot\Y_{sss}\X_s +3\Y_{ss}\cdot\Y_{sss}\bm{W}_s\big)\,ds\\
  &\quad +\int_\T \varphi\bigg(\X_{ss}\cdot\bigg(T_{m_\epsilon^{\rm n}}-\frac{\abs{\log\epsilon}}{4\pi}\bigg)+ \frac{\abs{\log\epsilon}}{4\pi}\bm{W}_{ss}\cdot \bigg)[\Y_{ssss}+3\Y_{ss}\cdot\Y_{sss}\Y_s] \,ds\,.
\end{aligned}
\end{equation}
%
%
Again using Lemma \ref{lem:mult_RFT_diff} and the inequalities \eqref{eq:interps}, we may estimate 
\begin{equation}
\begin{aligned}
  \abs{R^\perp} &\le c\norm{T_{m_\epsilon^{\rm n}}^{1/2}(\varphi\X_{ss})}_{L^2}\bigg(\norm{T_{m_\epsilon^{\rm n}}^{1/2}\bm{W}_{ssss}}_{L^2}+\norm{T_{m_\epsilon^{\rm n}}^{1/2}(\p_s\abs{\bm{W}_{ss}}^2\X_s)}_{L^2}\\
  &\qquad\qquad  +\norm{T_{m_\epsilon^{\rm n}}^{1/2}(\p_s(\bm{W}_{ss}\cdot\Y_{ss})\X_s)}_{L^2} +\norm{T_{m_\epsilon^{\rm n}}^{1/2}(\p_s\abs{\Y_{ss}}^2\bm{W}_s)}_{L^2}\bigg) \\
  &\qquad +c\bigg(\norm{T_{m_\epsilon^{\rm n}}^{1/2}(\varphi\X_{ss})}_{L^2}+ \abs{\log\epsilon}\norm{\varphi}_{L^2}\norm{\bm{W}_{ss}}_{L^\infty}\bigg)\big( \norm{\Y_s}_{H^3}+ \norm{\Y_{ss}}_{H^1}^2\big) \\
  %
  %
  &\le c\norm{T_{m_\epsilon^{\rm n}}^{1/2}(\varphi\X_{ss})}_{L^2}\bigg(\norm{T_{m_\epsilon^{\rm n}}^{1/2}\bm{W}_{ssss}}_{L^2}(1+\norm{\bm{W}_s}_{H^1}+\norm{\Y_s}_{H^3}) \\
  &\qquad\qquad +(1+\abs{\log\epsilon}^{1/2}\norm{\bm{W}_s}_{H^1})(\norm{\Y_s}_{H^3}+\norm{\Y_{ss}}_{H^1}^2) \bigg) \\
  &\qquad + c\abs{\log\epsilon}\norm{T_{m_\epsilon^{\rm n}}^{-1}T_{m_\epsilon^{\rm n}}(\varphi\X_s)\cdot\X_s}_{L^2}\norm{\bm{W}_{sss}}_{\dot H^{1/2}}^{2/3}\norm{\bm{W}_s}_{H^1}^{1/3}\big( \norm{\Y_s}_{H^3}+ \norm{\Y_{ss}}_{H^1}^2\big)\\
  %
  %
   &\le c\norm{T_{m_\epsilon^{\rm n}}^{1/2}(\varphi\X_{ss})}_{L^2}\bigg(\norm{T_{m_\epsilon^{\rm n}}^{1/2}\bm{W}_{ssss}}_{L^2}(1+\norm{\bm{W}_s}_{H^1}+\norm{\Y_s}_{H^3}) \\
  &\qquad\qquad +(1+\abs{\log\epsilon}^{1/2}\norm{\bm{W}_s}_{H^1})(\norm{\Y_s}_{H^3}+\norm{\Y_{ss}}_{H^1}^2) \bigg) \\
  &\qquad + c\abs{\log\epsilon}^{1/2}\norm{T_{m_\epsilon^{\rm n}}^{1/2}(\varphi\X_s)_s}_{L^2}\norm{\bm{W}_{sss}}_{\dot H^{1/2}}^{2/3}\norm{\bm{W}_s}_{H^1}^{1/3}\big( \norm{\Y_s}_{H^3}+ \norm{\Y_{ss}}_{H^1}^2\big)\,.
\end{aligned}
\end{equation}
Here we have reused some of the estimates from the LHS bound~\eqref{eq:LHS_diffbd} and the $R^\parallel$ bound~\eqref{eq:Rpar_bd}.



Altogether, taking $\varphi=\tau_W$ in \eqref{eq:diffXY_weak} and combining the bounds for LHS and RHS, we may estimate 
\begin{equation}\label{eq:BXtautau}
\begin{aligned}
  &\abs{\bm{B}_\epsilon^X(\tau_W,\tau_W)}\le \abs{\rm LHS} + \abs{\rm RHS}\\
  %
  &\le c\bigg( \norm{T_{m_\epsilon^{\rm t}}^{1/2}((\tau_W)_s\X_s)}_{L^2} + 
  \norm{T_{m_\epsilon^{\rm n}}^{1/2}(\tau_W\X_{ss})}_{L^2}\bigg)\bigg(\abs{\log\epsilon}^{1/2}\norm{\bm{W}_s}_{H^1}(\norm{\Y_s}_{H^3}+\norm{\Y_s}_{H^2}^2)\\
  &\qquad  + \abs{\log\epsilon}^{1/2}\norm{\bm{W}_s}_{H^1}^{4/3}\norm{\bm{W}_{sss}}_{\dot H^{1/2}}^{2/3}(\norm{\Y_s}_{H^3}+\norm{\Y_{ss}}_{H^1}^2)\\
  &\qquad\quad + \norm{T_{m_\epsilon^{\rm t}}^{1/2}\bm{W}_{ssss}}_{L^2}(1+\norm{\bm{W}_s}_{H^1}+\norm{\Y_s}_{H^3})\\
  &\qquad\qquad + (\norm{\Y_s}_{H^3}+\norm{\Y_{ss}}_{H^1}^2)\big(1 + \norm{\X_s}_{H^1}^2 + \norm{\Y_{ss}}_{H^1}^2\big)\bigg)\\
  %
  &\le c\bigg( \norm{T_{m_\epsilon^{\rm t}}^{1/2}((\tau_W)_s\X_s)}_{L^2} + 
  \norm{T_{m_\epsilon^{\rm n}}^{1/2}(\tau_W\X_{ss})}_{L^2}\bigg)\bigg(\abs{\log\epsilon}^{1/2}\norm{\bm{W}_s}_{H^1}(\norm{\Y_s}_{H^3}+\norm{\Y_s}_{H^2}^2)\\
  &\qquad\qquad  + \norm{T_{m_\epsilon^{\rm t}}^{1/2}\bm{W}_{ssss}}_{L^2}(1+\norm{\X_s}_{H^1}+\norm{\Y_s}_{H^3}+\norm{\Y_{ss}}_{H^1}^2)\\
  &\qquad\qquad\quad + (\norm{\Y_s}_{H^3}+\norm{\Y_{ss}}_{H^1}^2)\big(1 + \norm{\X_s}_{H^1}^2 + \norm{\Y_{ss}}_{H^1}^2\big)\bigg)\,.
\end{aligned}
\end{equation}
Here, to simplify, we have used that $m_\epsilon^{\rm t}(k)\lesssim m_\epsilon^{\rm n}(k) \lesssim m_\epsilon^{\rm t}(k)$, so we may interchange $T_{m_\epsilon^{\rm t}}$ and $T_{m_\epsilon^{\rm n}}$. In addition, we have replaced one instance of $\norm{\bm{W}_s}_{H^1}\le \norm{\X_s}_{H^1}+\norm{\Y_s}_{H^1}$, and used that
\begin{equation}
\begin{aligned}
  \norm{\bm{W}_{sss}}_{\dot H^{1/2}}^{2/3}\norm{\bm{W}_s}_{H^1}^{1/3}&\le c\abs{\log\epsilon}^{-1/3}\norm{T_{m_\epsilon^{\rm t}}\bm{W}_{ssss}}_{L^2}^{2/3}\norm{\bm{W}_s}_{H^1}^{1/3}\\
  &\le c(\abs{\log\epsilon}^{-1/2}\norm{T_{m_\epsilon^{\rm t}}\bm{W}_{ssss}}_{L^2} + \norm{\bm{W}_s}_{H^1})\,.
\end{aligned}
\end{equation} 
%
%
Finally, using that 
\begin{equation}
  \bm{B}_\epsilon^X(\tau_W,\tau_W) = \norm{T_{m_\epsilon^{\rm t}}^{1/2}((\tau_W)_s\X_s)}_{L^2}^2 + 
  \norm{T_{m_\epsilon^{\rm n}}^{1/2}(\tau_W\X_{ss})}_{L^2}^2\,,
\end{equation}
applying Young's inequality to \eqref{eq:BXtautau} then yields \eqref{eq:tauWbd}.
%
%
\end{proof}

\end{document}
